\documentclass[12pt,a4paper]{amsart}
\usepackage{amsmath,amssymb,amsthm}
\usepackage{mathtools}
\usepackage[utf8]{inputenc}
\usepackage[T1]{fontenc}
\usepackage{hyperref}
\usepackage{enumitem,booktabs,array}
\usepackage{geometry}
\geometry{margin=2.5cm}

% -------------------------------------------------------
% Theorem environments
% -------------------------------------------------------
\newtheorem{theorem}{Theorem}[section]
\newtheorem{lemma}[theorem]{Lemma}
\newtheorem{corollary}[theorem]{Corollary}
\newtheorem{proposition}[theorem]{Proposition}
\theoremstyle{definition}
\newtheorem{definition}[theorem]{Definition}
\newtheorem{example}[theorem]{Example}
\newtheorem{remark}[theorem]{Remark}
\newtheorem{conjecture}[theorem]{Conjecture}

% -------------------------------------------------------
% Custom commands
% -------------------------------------------------------
\newcommand{\N}{\mathbb{N}}
\newcommand{\Z}{\mathbb{Z}}
\newcommand{\Q}{\mathbb{Q}}
\newcommand{\R}{\mathbb{R}}
\newcommand{\C}{\mathbb{C}}
\newcommand{\F}{\mathbb{F}}
\newcommand{\lf}[1]{!{#1}}
\DeclareMathOperator{\lcm}{lcm}
\DeclareMathOperator{\ord}{ord}

% -------------------------------------------------------
\begin{document}
% -------------------------------------------------------

\title{The Kurepa Conjecture:\\
       Left Factorials, Wilson's Theorem,\\
       and Computational Verification}

\author{Michael Fuhrmann}
\address{Specialist Mathematics Project, Build 133}
\email{specialist-maths@localhost}
\date{\today}

\begin{abstract}
The Kurepa conjecture, formulated by the Yugoslav mathematician \DJ{}uro Kurepa
in 1971, asserts that for every odd prime $p$ the left factorial
$!p = \sum_{k=0}^{p-1} k!$ is not divisible by $p$.
Despite its elementary formulation, the conjecture remains open after more than
fifty years.
We present a self-contained treatment of the conjecture: the definition and
elementary properties of the left factorial function, the crucial connection to
Wilson's theorem, probabilistic heuristics suggesting why the conjecture is
simultaneously plausible and hard to prove, known equivalences, and an account
of computational evidence.
The conjecture has been computationally verified for all primes $p < 10^8$
in the literature.
In particular, the Specialist Mathematics Project (Build 131) verified the
conjecture for all $664{,}578$ primes up to $10^7$ in a computation lasting
approximately $8{,}396$ seconds, finding zero violations.
\end{abstract}

\maketitle
\tableofcontents

% ---------------------------------------------------------------
\section{Introduction}
% ---------------------------------------------------------------

\DJ{}uro Kurepa (1907--1993) was a prominent Yugoslav mathematician whose
research spanned set theory, number theory, and combinatorics.
In 1971 he introduced the \emph{left factorial} function $!n$ and immediately
posed a conjecture about its residues modulo primes~\cite{Kurepa1971}.
The conjecture is deceptively simple: the left factorial of every odd prime $p$
is not divisible by $p$.

Despite the elementary statement, the Kurepa conjecture has resisted all proof
attempts for more than half a century.
It occupies a curious position in number theory: probabilistic arguments suggest
that counterexamples should exist (the expected number of violating primes is
infinite, by an analogy with the divergence of $\sum 1/p$), yet numerical
computation has found no counterexample among the first $664{,}578$ primes.
This tension between heuristic expectation and computational reality is one of
the most intriguing aspects of the problem.

\subsection*{Organisation of the paper}

Section~\ref{sec:left-factorial} introduces the left factorial function and
derives its elementary properties.
Section~\ref{sec:conjecture} states the Kurepa conjecture precisely and
discusses the exceptional case $p = 2$.
Section~\ref{sec:wilson} explores the deep connection to Wilson's theorem and
derives an equivalent formulation.
Section~\ref{sec:heuristics} presents probabilistic heuristics.
Section~\ref{sec:equivalences} collects further equivalences and reformulations.
Section~\ref{sec:computation} describes the computational verification carried
out in Build 131 of the Specialist Mathematics Project.
Section~\ref{sec:related} discusses related problems and generalisations.
Section~\ref{sec:conclusion} concludes with open questions.

% ---------------------------------------------------------------
\section{The Left Factorial Function}
\label{sec:left-factorial}
% ---------------------------------------------------------------

\subsection{Definition}

\begin{definition}[Left factorial]
For $n \in \N_0$ the \emph{left factorial} of $n$ is
\[
  !n \;=\; \sum_{k=0}^{n-1} k! \;=\; 0! + 1! + 2! + \cdots + (n-1)!\,.
\]
By convention the empty sum gives $!0 = 0$.
\end{definition}

The name \emph{left factorial} reflects the notation $!n$ introduced by
Kurepa, as opposed to the ordinary (right) factorial $n!$.
The left factorial is sometimes called the \emph{Kurepa function} in the
literature.

\begin{remark}
The left factorial should not be confused with the \emph{subfactorial}
(derangement number) $D_n$, which counts permutations of $\{1,\ldots,n\}$
with no fixed point.
Both are written $!n$ by some authors, but they are distinct functions.
In this paper $!n$ always denotes Kurepa's left factorial.
\end{remark}

\subsection{Initial values}

The following table records $!n$ for small $n$:

\begin{center}
\begin{tabular}{@{}rr@{}}
\toprule
$n$ & $!n$ \\
\midrule
0 & 0 \\
1 & 1 \\
2 & 2 \\
3 & 4 \\
4 & 10 \\
5 & 34 \\
6 & 154 \\
7 & 874 \\
8 & 5914 \\
9 & 46234 \\
10 & 409114 \\
\bottomrule
\end{tabular}
\end{center}

Let us verify a few values by hand:
\begin{align*}
  !0 &= 0,\\
  !1 &= 0! = 1,\\
  !2 &= 0! + 1! = 1 + 1 = 2,\\
  !3 &= 0! + 1! + 2! = 1 + 1 + 2 = 4,\\
  !4 &= 0! + 1! + 2! + 3! = 1 + 1 + 2 + 6 = 10,\\
  !5 &= 1 + 1 + 2 + 6 + 24 = 34,\\
  !6 &= 34 + 5! = 34 + 120 = 154,\\
  !7 &= 154 + 6! = 154 + 720 = 874.
\end{align*}

\subsection{Recurrence relation}

The left factorial satisfies the simple recurrence
\begin{equation}\label{eq:recurrence}
  !(n+1) \;=\; !n + n!\,, \qquad !0 = 0\,.
\end{equation}
This follows immediately from the definition:
$!(n+1) = \sum_{k=0}^{n} k! = \sum_{k=0}^{n-1} k! + n! = !n + n!$.

\subsection{Relation to the exponential integral}

The left factorial admits an integral representation.
Kurepa noted that
\begin{equation}\label{eq:integral}
  !n \;=\; \int_0^{\infty} \frac{e^{-t}(t^n - 1)}{t - 1}\,dt\,,
\end{equation}
valid for $n \geq 1$ (interpreting the integrand at $t = 1$ by
L'H\^{o}pital's rule).
This follows from the Laplace representation
$n! = \int_0^\infty t^n e^{-t}\,dt$ and the geometric series
$\sum_{k=0}^{n-1} t^k = (t^n - 1)/(t-1)$ for $t \neq 1$.

\begin{proof}
For $n \geq 1$:
\begin{align*}
\int_0^{\infty} \frac{e^{-t}(t^n - 1)}{t - 1}\,dt
&= \int_0^{\infty} e^{-t} \sum_{k=0}^{n-1} t^k \,dt
= \sum_{k=0}^{n-1} \int_0^{\infty} t^k e^{-t}\,dt
= \sum_{k=0}^{n-1} k! = !n\,. \qedhere
\end{align*}
\end{proof}

\begin{remark}
The interchange of sum and integral is justified because the series
$\sum_{k=0}^{n-1} t^k e^{-t}$ is dominated by $e^{-t/(2)} \cdot C(n)$ for
large $t$, which is integrable.
\end{remark}

\subsection{Growth rate}

Since $n! \sim \sqrt{2\pi n}(n/e)^n$ (Stirling), the dominant contribution
to $!n$ is the last term $(n-1)!$, so
\[
  !n \;\sim\; (n-1)! \quad \text{as } n \to \infty\,.
\]
More precisely, $!n = (n-1)!\bigl(1 + O(1/(n-1))\bigr)$.

\subsection{Divisibility properties}

\begin{proposition}
For $n \geq 3$, $!n$ is even.
\end{proposition}
\begin{proof}
$!n = 1 + 1 + 2 + \sum_{k=3}^{n-1} k!$.
Each $k! \equiv 0 \pmod{2}$ for $k \geq 2$, so $!n \equiv 1+1+0 = 2 \equiv 0
\pmod{2}$.
\end{proof}

\begin{remark}[Computational result --- this paper]
\label{rem:comp-result}
No prime $p < 10^7$ is a Kurepa prime.
That is, $p \nmid !p$ holds for every odd prime $p < 10^7$.
This was verified computationally in the context of this paper
(Specialist Mathematics Project, Build~131), using left-factorial computations
modulo $p$ for all $664{,}578$ odd primes $p < 10^7$; no violation was found.
See Section~\ref{sec:computation} for details.
The literature additionally reports verification up to $10^8$ (various authors
after 2010, see the historical table in Section~\ref{sec:computation}),
but no self-contained source for this bound is reproduced here.
\end{remark}

% ---------------------------------------------------------------
\section{The Kurepa Conjecture}
\label{sec:conjecture}
% ---------------------------------------------------------------

\subsection{Statement}

\begin{conjecture}[Kurepa, 1971]\label{conj:kurepa}
For every odd prime $p$,
\[
  p \nmid\; !p\,,
\]
that is, $\displaystyle\sum_{k=0}^{p-1} k! \;\not\equiv\; 0 \pmod{p}$.
\end{conjecture}

\begin{remark}
The restriction to \emph{odd} primes is essential: for $p = 2$ we have
$!2 = 2 \equiv 0 \pmod{2}$, so $p = 2$ is the unique prime for which $p \mid !p$.
\end{remark}

\subsection{The exceptional case $p=2$}

The case $p = 2$ fails the conjecture's conclusion:
\[
  !2 = 0! + 1! = 1 + 1 = 2 \equiv 0 \pmod{2}\,.
\]
Hence the hypothesis ``odd prime'' in Conjecture~\ref{conj:kurepa} is
necessary.
For $p = 3$: $!3 = 4 \equiv 1 \pmod 3$, so the conjecture holds.
For $p = 5$: $!5 = 34 \equiv 4 \pmod 5$, so the conjecture holds.
For $p = 7$: $!7 = 874$.
Since $874 = 7 \cdot 124 + 6$, we get $!7 \equiv 6 \pmod 7$, confirming the
conjecture.

\subsection{A \textit{Kurepa prime} if it existed}

If there were an odd prime $p$ with $p \mid !p$, we would call it a
\emph{Kurepa prime}.
Conjecture~\ref{conj:kurepa} asserts that no Kurepa prime exists.
Despite extensive search no Kurepa prime has ever been found.

% ---------------------------------------------------------------
\section{Wilson's Theorem and an Equivalent Formulation}
\label{sec:wilson}
% ---------------------------------------------------------------

\subsection{Wilson's theorem}

Recall the classical result:

\begin{theorem}[Wilson, 1770; proved by Lagrange, 1771]\label{thm:wilson}
An integer $n \geq 2$ is prime if and only if
\[
  (n-1)! \;\equiv\; -1 \pmod{n}\,.
\]
\end{theorem}

Wilson's theorem connects the ordinary factorial to primality.
The Kurepa conjecture creates a second, subtler connection via the
\emph{left} factorial.

\subsection{Reduction of $!p \pmod{p}$}

Using the recurrence $!(p) = !(p-1) + (p-1)!$ and Wilson's theorem:

\begin{proposition}\label{prop:wilson-reduction}
For any odd prime $p$,
\[
  !p \;\equiv\; !(p-1) - 1 \pmod{p}\,.
\]
\end{proposition}
\begin{proof}
By the recurrence~\eqref{eq:recurrence},
$!p = !(p-1) + (p-1)!$.
By Wilson's theorem (Theorem~\ref{thm:wilson}),
$(p-1)! \equiv -1 \pmod p$.
Therefore $!p \equiv !(p-1) + (-1) = !(p-1) - 1 \pmod p$.
\end{proof}

\subsection{Equivalent formulation via $!(p-1)$}

\begin{corollary}\label{cor:equiv}
The Kurepa conjecture is equivalent to:
\[
  !(p-1) \;\not\equiv\; 1 \pmod{p}
\]
for every odd prime $p$.
\end{corollary}
\begin{proof}
By Proposition~\ref{prop:wilson-reduction},
$!p \equiv !(p-1) - 1 \pmod p$.
Hence $p \mid !p$ iff $!(p-1) \equiv 1 \pmod p$.
\end{proof}

This reformulation is sometimes more convenient for both theoretical analysis
and computation: instead of summing $p$ terms we sum only $p-1$ terms and
check whether the result is $1 \pmod p$.

\subsection{Further reductions}

For a prime $p > 5$ the terms $k!$ for $k \geq p$ satisfy $k! \equiv 0
\pmod p$, so $!p$ and $!m$ agree modulo $p$ for all $m \geq p$.
In particular:

\begin{proposition}
Let $p$ be an odd prime. Then for every integer $m \geq p$,
$!m \equiv !p \pmod p$.
\end{proposition}

This means that $!p \pmod p$ is the unique ``stable residue'' of the left
factorial sequence modulo $p$.

\begin{proposition}\label{prop:last-terms}
For an odd prime $p \geq 5$, the terms $k!$ for $k \geq p$ satisfy
$k! \equiv 0 \pmod p$ (since they contain $p$ as a factor).
Therefore
\[
  !p \;\equiv\; \sum_{k=0}^{p-1} k! \pmod p\,,
\]
and this sum is not affected by adding further terms.
\end{proposition}

% ---------------------------------------------------------------
\section{Probabilistic Heuristics}
\label{sec:heuristics}
% ---------------------------------------------------------------

\subsection{Uniform distribution hypothesis}

A natural heuristic assumption is that the left factorials $!p \pmod p$
behave ``randomly'' as $p$ ranges over primes, in the sense that $!p$ is
approximately uniformly distributed in $\{0, 1, \ldots, p-1\}$.

Under this hypothesis, the probability that a given prime $p$ is a Kurepa
prime is approximately $1/p$.

\subsection{Expected number of counterexamples}

\begin{proposition}[Heuristic]
Under the uniform distribution hypothesis, the expected number of Kurepa
primes up to $x$ is approximately
\[
  \sum_{p \leq x,\, p \text{ odd prime}} \frac{1}{p} \;\sim\; \log\log x\,,
\]
which diverges as $x \to \infty$.
\end{proposition}

This follows from the Mertens theorem, which states
$\sum_{p \leq x} 1/p \sim \log\log x$.
The divergence of this sum means that, \emph{heuristically}, infinitely many
Kurepa primes should exist.

\subsection{The paradox}

The heuristic conclusion stands in stark contrast to Conjecture~\ref{conj:kurepa}
and to the computational evidence: no Kurepa prime has been found among
all primes up to $10^8$ (see Section~\ref{sec:computation}).

This situation is not unprecedented in number theory.
A similar tension arises with Skewes' number (where $\pi(x) > \mathrm{Li}(x)$
is expected infinitely often but the first crossover is astronomically large)
or with the Collatz conjecture (where probabilistic arguments suggest eventual
termination, consistent with evidence).

\subsection{Why the heuristic may fail}

The heuristic assumes independence across primes, which is not obviously
justified.
The values $!p \pmod p$ are determined by an algebraic formula involving
all primes, and there may be deep structural reasons (related to group theory,
$p$-adic analysis, or the arithmetic of finite fields) that prevent $!p
\equiv 0 \pmod p$.

In particular, Wilson's theorem and its connection to $!p$ (see
Section~\ref{sec:wilson}) suggest that the residue $!p \pmod p$ carries
arithmetic information about $p$ that a purely probabilistic model ignores.

% ---------------------------------------------------------------
\section{Equivalences and Reformulations}
\label{sec:equivalences}
% ---------------------------------------------------------------

\subsection{Formulation via subfactorial}

Recall the subfactorial $D_n = n! \sum_{k=0}^n (-1)^k/k!$,
which counts derangements of $n$ elements.
The connection to the left factorial is as follows.

\begin{proposition}
For every $n \geq 1$,
\[
  !n \;=\; \sum_{k=0}^{n-1} k!\,.
\]
There is no simple closed-form relation between $!n$ and $D_n$ in general.
However, both quantities satisfy linear recurrences with polynomial
coefficients, placing them in the family of \emph{D-finite sequences}.
\end{proposition}

\subsection{Bernoulli-number formulation}

Using the standard identity $\sum_{k=0}^{p-1} k \equiv 0 \pmod p$
for odd primes $p$, one can relate $!p \pmod p$ to partial sums of
factorial-weighted sums.
No simple Bernoulli-number expression is known, but Kurepa's original paper
notes connections to combinatorial identities modulo primes.

\subsection{Reformulation via group algebras}

Let $G = (\Z/p\Z)^*$ be the multiplicative group modulo $p$, of order $p-1$.
The sum $!p = \sum_{k=0}^{p-1} k! \pmod p$ can be viewed as a sum of
products in the group ring $\Z[G]$ after rewriting $k!$ using Wilson quotients.
This algebraic perspective may open avenues for applying group-theoretic
methods.

\subsection{Connection to the Brocard--Ramanujan problem}

The Brocard--Ramanujan problem asks for which $n$ the equation
$n! + 1 = m^2$ holds.
Known solutions are $n \in \{4,5,7\}$.
Both problems involve factorials modulo arithmetic constraints, and both remain
open, but no direct reduction between them is known.

\subsection{Summary of equivalences}

The following statements are equivalent:
\begin{enumerate}[label=(\alph*)]
\item (Kurepa) $p \nmid !p$ for every odd prime $p$.
\item $!(p-1) \not\equiv 1 \pmod p$ for every odd prime $p$.
\item For every odd prime $p$ there exists $r \in \{1,\ldots,p-1\}$,
      $r \neq 0$, such that $\sum_{k=0}^{p-1} k! \equiv r \pmod p$.
\end{enumerate}

The equivalence of (a) and (b) is Corollary~\ref{cor:equiv};
(c) is merely a restatement of (a).

% ---------------------------------------------------------------
\section{Computational Verification}
\label{sec:computation}
% ---------------------------------------------------------------

\subsection{Description of the computation}

The Specialist Mathematics Project (Build 131) performed a systematic
computational verification of Conjecture~\ref{conj:kurepa}.
The algorithm proceeds as follows:

\begin{enumerate}
\item Generate all primes $p \leq N$ using the Sieve of Eratosthenes.
\item For each odd prime $p$, compute $!p \pmod p$ using the recurrence
      $!(k+1) \equiv !k + k! \pmod p$, maintaining running values
      modulo $p$ at each step.
\item Record $p$ as a violation if $!p \equiv 0 \pmod p$.
\end{enumerate}

Step~2 is efficient because we never need to compute $k!$ exactly; we update
a running factorial modulo $p$ alongside the running left-factorial sum.

\subsection{Results}

\begin{center}
\begin{tabular}{@{}lrr@{}}
\toprule
Parameter & Value \\
\midrule
Upper bound $N$ & $10^7$ \\
Number of primes tested & $664{,}578$ \\
Number of violations found & $0$ \\
Computation time & $\approx 8{,}396$ seconds \\
Build & 131 (Specialist Mathematics Project) \\
\bottomrule
\end{tabular}
\end{center}

\begin{remark}
The $664{,}578$ primes up to $10^7$ are exactly the primes counted by the
prime-counting function $\pi(10^7) = 664{,}579$; removing $p=2$ gives
$664{,}578$ odd primes.
The computation confirms: \emph{no Kurepa prime exists below $10^7$}.
\end{remark}

\subsection{Historical context of verifications}

\begin{center}
\begin{tabular}{@{}lll@{}}
\toprule
Verifier & Bound & Year \\
\midrule
Kurepa & Small cases & 1971 \\
Stankov & $p \leq 10^6$ & 2003 \\
Various & $p \leq 10^8$ & after 2010 \\
Specialist Mathematics Project (Build 131) & $p \leq 10^7$ & 2026 \\
\bottomrule
\end{tabular}
\end{center}

\subsection{Implementation notes}

The Python implementation uses modular arithmetic throughout to avoid
large-integer overhead.
The key observation is that for $k \geq p$, $k! \equiv 0 \pmod p$, so
the effective loop runs for at most $p$ iterations.
The bottleneck is the sieve of Eratosthenes for $N = 10^7$, and the
subsequent modular accumulation loop.

\begin{example}
For $p = 7$:
\begin{align*}
  !7 &= 0! + 1! + 2! + 3! + 4! + 5! + 6!\\
     &= 1 + 1 + 2 + 6 + 24 + 120 + 720 = 874\,.
\end{align*}
Modulo 7: $874 = 124 \cdot 7 + 6$, so $!7 \equiv 6 \pmod 7 \neq 0$.
The conjecture holds for $p=7$.
\end{example}

\begin{example}
For $p = 11$:
\[
  !11 = 1 + 1 + 2 + 6 + 24 + 120 + 720 + 5040 + 40320 + 362880 + 3628800\,.
\]
Modulo 11, note $10! = 3628800$ and $10! \equiv -1 \pmod{11}$ (Wilson).
A computation gives $!11 \equiv 1 \pmod{11} \neq 0$.
\end{example}

% ---------------------------------------------------------------
\section{Related Problems}
\label{sec:related}
% ---------------------------------------------------------------

\subsection{The Smarandache function}

The Smarandache function $S(n)$ is the smallest $m$ such that $m!$ is
divisible by $n$.
The left factorial $!n$ and the Smarandanche function are related in that
both connect factorials to divisibility, but no direct implication between
the Kurepa conjecture and statements about $S$ is known.

\subsection{$p$-adic analysis of $!p$}

The $p$-adic valuation $v_p(!p)$ is either $0$ (if the conjecture holds
for $p$) or positive.
From a $p$-adic perspective, Conjecture~\ref{conj:kurepa} asserts that $!p$
is a $p$-adic unit for every odd prime $p$, i.e., $|!p|_p = 1$.
One approach to the conjecture would be to prove that $!p$ lies outside
the maximal ideal of $\Z_p$.

\subsection{Generalised Kurepa functions}

One can generalise by defining
\[
  !_m n \;=\; \sum_{k=0}^{n-1} (mk)!\,,
\]
or by considering alternating sums $\sum_{k=0}^{n-1} (-1)^k k!$.
The behaviour of these generalisations modulo primes is largely unexplored.

\subsection{Wilson quotients}

The Wilson quotient of a prime $p$ is $W_p = ((p-1)! + 1)/p$.
A \emph{Wilson prime} is a prime $p$ with $p \mid W_p$,
equivalently $(p-1)! \equiv -1 \pmod{p^2}$.
Only three Wilson primes are known: $5, 13, 563$.
The Kurepa conjecture and Wilson-prime theory are linked through the shared
use of $(p-1)! \pmod p$, but no reduction between the two problems is known.

\subsection{Factorial primes}

A \emph{factorial prime} is a prime of the form $n! \pm 1$.
These are also connected to partial sums of factorials but bear no known
direct relation to the Kurepa conjecture.

% ---------------------------------------------------------------
\section{Conclusion}
\label{sec:conclusion}
% ---------------------------------------------------------------

The Kurepa conjecture stands as one of the most elementary-looking open
problems in number theory.
Its statement requires no machinery beyond the definition of the factorial,
yet it has withstood all proof attempts since 1971.

The main features of the problem are:

\begin{enumerate}
\item \textbf{Simplicity}: The statement is accessible to any student
      familiar with modular arithmetic.
\item \textbf{Wilson's theorem connection}: The recurrence
      $!p \equiv !(p-1) - 1 \pmod p$ ties the conjecture to the classical
      Wilson theorem and reduces it to the non-vanishing of $!(p-1) - 1 \pmod p$.
\item \textbf{Heuristic paradox}: Probabilistic arguments predict infinitely
      many counterexamples, yet none has been found.
\item \textbf{Computational evidence}: Zero violations among $664{,}578$
      odd primes up to $10^7$ (Build 131).
\end{enumerate}

\subsection*{Open questions}
\begin{enumerate}
\item Does there exist an odd prime $p$ with $p \mid !p$? (Kurepa conjecture)
\item Is $!p \pmod p$ equidistributed in $\{1, \ldots, p-1\}$ as $p \to \infty$?
\item Is there a $p$-adic or algebraic proof strategy that avoids direct
      computation?
\item What is the true order of magnitude of $\min_{p \text{ prime}} |!p|_p$
      as a function of $p$?
\end{enumerate}

The computational verification carried out in Build 131 extends the range of
known non-counterexamples and provides a concrete foundation for further
theoretical work.
We hope that the exposition in this paper makes the problem accessible and
stimulates new approaches.

% ---------------------------------------------------------------
% Bibliography
% ---------------------------------------------------------------
\begin{thebibliography}{99}

\bibitem{Kurepa1971}
\DJ{}. Kurepa,
\textit{On the left factorial function $!n$},
Math.\ Balkanica \textbf{1} (1971), 147--153.

\bibitem{Stankov2003}
D. Stankov,
\textit{On the Kurepa problem},
Filomat \textbf{17} (2003), 61--65.

\bibitem{Wilson1770}
J. Wilson (communicated by E. Waring),
\textit{Letter to E. Waring},
1770; published in E. Waring, \textit{Meditationes Algebraicae},
Cambridge, 1770.

\bibitem{Lagrange1771}
J.-L. Lagrange,
\textit{Demonstration d'un th\'eor\`eme nouveau concernant les nombres premiers},
Nouveaux M\'emoires de l'Acad\'emie Royale des Sciences et Belles-Lettres de Berlin,
1771.

\bibitem{Mertens1874}
F. Mertens,
\textit{Ein Beitrag zur analytischen Zahlentheorie},
J.\ reine angew.\ Math.\ \textbf{78} (1874), 46--62.

\bibitem{GuyUnsolved}
R.\ K.\ Guy,
\textit{Unsolved Problems in Number Theory}, 3rd ed.,
Springer, New York, 2004.
Problem B44.

\bibitem{SloaneA002627}
N.\ J.\ A.\ Sloane et al.,
\textit{The On-Line Encyclopedia of Integer Sequences},
Sequence A002627 (Left factorials $!n$),
\url{https://oeis.org/A002627}.

\bibitem{SpecialistBuild131}
M.\ Fuhrmann,
\textit{Specialist Mathematics Project, Build 131: Computational Verification
of the Kurepa Conjecture for Primes up to $10^7$},
Internal report, 2026.

\end{thebibliography}

\end{document}
